\begin{table*}[!t]
\centering
\footnotesize
\setlength{\tabcolsep}{4pt}
\caption{ChromaDB (HNSW, busca aproximada) e FAISS (\texttt{IndexFlatL2}, busca exata) sobre os mesmos vetores BGE-M3. Top-4 idêntico: consultas em que os dois bancos entregam os mesmos quatro fragmentos. Na última coluna, a métrica principal do Stuff do gpt-4o-mini nas 100 perguntas, refazendo com o top-4 do FAISS as perguntas em que ele difere do índice dos experimentos.}
\label{tab:bancos}
\begin{tabular}{lrcccccccc}
\toprule
 & & \multicolumn{2}{c}{Indexação (s)} & \multicolumn{2}{c}{Consulta (ms)} & \multicolumn{2}{c}{Memória (MB)} & & \\
\cmidrule(lr){3-4}\cmidrule(lr){5-6}\cmidrule(lr){7-8}
Dataset & Vetores & Chroma & FAISS & Chroma & FAISS & Chroma & FAISS & Top-4 idêntico & Stuff: Chroma × FAISS \\
\midrule
FaQuAD & 318 & 0,63 & 0,001 & 3,78 & 0,05 & 29,3 & 1,6 & 100\% & 0,86 × 0,86 \\
Docugami SEC 10-Q & 4.463 & 4,01 & 0,017 & 2,11 & 2,08 & 191,9 & 17,8 & 100\% & 0,09 × 0,09 \\
MedQA & 121.553 & 260,83 & 0,304 & 5,00 & 36,47 & 801,5 & 474,8 & 82\% & 0,68 × 0,69 \\
\bottomrule
\end{tabular}
\end{table*}
